% Generated draft: requires mathematical and source review.
\documentclass[11pt]{article}
\usepackage[margin=25mm]{geometry}
\usepackage{fontspec}
\setmainfont{Latin Modern Roman}
\usepackage{amsmath,amssymb,mathtools,unicode-math}
\setmathfont{Latin Modern Math}
\usepackage{enumitem,needspace,xurl,hyperref}
\setlength{\parindent}{0pt}
\setlength{\parskip}{5pt}
\setlength{\emergencystretch}{3em}
\newcommand{\N}{\mathbb{N}}
\newcommand{\Z}{\mathbb{Z}}
\newcommand{\Q}{\mathbb{Q}}
\newcommand{\R}{\mathbb{R}}
\newcommand{\C}{\mathbb{C}}
\newcommand{\F}{\mathbb{F}}
\newcommand{\A}{\mathbb{A}}
\newcommand{\PP}{\mathbb{P}}
\newcommand{\E}{\mathbb{E}}
\newcommand{\eps}{\varepsilon}
\newcommand{\dd}{\,\mathrm{d}}
\begin{document}
\section*{UnsolvedMath ID: 700924}
\textbf{Source problem code:} catalogue entry
\subsection{Definitions and mathematical statement}
A group \(G\) is locally finite if every finitely generated subgroup of \(G\) is finite. For \(g\in G\), its centralizer is
\[
C_G(g)=\{x\in G:xg=gx\}.
\]
For a prime \(p\), the quasicyclic \(p\)-group is
\[
C_{p^\infty}=\{z\in\C^\times:z^{p^n}=1\text{ for some integer }n\geq 0\},
\]
with multiplication inherited from \(\C^\times\). A group \(K\) is Chernikov if it has a normal subgroup \(D\) of finite index such that
\[
D\cong\prod_{i=1}^{r}C_{p_i^\infty}
\]
for some integer \(r\geq 0\) and primes \(p_1,\ldots,p_r\), which need not be distinct. The empty product is the trivial group, so finite groups are included.

A group \(G\) is almost soluble if it has a soluble subgroup \(H\) of finite index. Here soluble means that the derived series, defined by \(H^{(0)}=H\) and \(H^{(n+1)}=[H^{(n)},H^{(n)}]\), satisfies \(H^{(d)}=\{1\}\) for some integer \(d\geq 0\).

The question is whether, for every locally finite group \(G\),
\[
\bigl(\exists g\in G\; C_G(g)\text{ is Chernikov}\bigr)
\quad\Longrightarrow\quad
\bigl(\exists H\leq G\; [G:H]<\infty\text{ and }H\text{ is soluble}\bigr).
\]
The supplied statement credits B. Hartley, whereas the proposer metadata names V. V. Belyaev; their respective roles are not explained in the record. The supplied registry labels the problem open; this is a source label, not an independently verified assessment.
\subsection{Short English statement}
Must every locally finite group containing an element whose centralizer is Chernikov have a soluble subgroup of finite index? The supplied record labels this question open.
\subsection{Sources}
\begin{itemize}
\item Local UnsolvedMath registry, supplied complete source record: permanent ID \(700924\), exact source problem code \texttt{catalogue entry}, upstream revision \texttt{d371ee917ac9647fbb06cd6fc529deec6f3dda82}. The record was updated on 2026-10-16.
\item \emph{The Kourovka Notebook (Unsolved Problems in Group Theory)}, Problem 13.6, 13th issue (1995); the supplied background also identifies the 21st issue edition (2026), arXiv:1401.0300v48, \url{https://arxiv.org/abs/1401.0300v48}. The supplied collection description names E. I. Khukhro and V. D. Mazurov as editors.
\item Attribution as supplied: proposer V. V. Belyaev, proposed year 1995; the statement itself begins with the attribution to B. Hartley. These attributions are retained without assigning unrecorded roles.
\end{itemize}
\end{document}
